\subsection{Carry the conversion through the whole company}
\label{sec:company-eps-join}

The EPS calculation must use earnings and shares from the same company
history. Return to the parent and subsidiary in
Section~\ref{sec:company-composition}. Their unchanged annual operations
produce group profit 102, of which7.35 belongs to non-controlling interests
and 94.65 to the parent's ordinary shareholders. Group cash is 644 and equity
724. The opening matured convertible was already redeemed in January;
the new instrument below is a separate July issue. Supply a register of ten
parent ordinary shares outstanding from the start of the year, with no
treasury shares, splits, preference claims or discontinued operations.
Capital 600 is a monetary balance and does not establish that share count.

On 1 July the parent receives 100 for a new instrument that matures on
1 October. Its original terms require a coupon of 4, followed by either
principal repayment of 92 or conversion of that principal into six ordinary
shares. The fixed conversion feature qualifies as equity in the parent's
functional currency; no fees, participation or other consequential earnings
effects are assumed. At the supplied comparable 92-day debt rate of $1/23$,
\[
 L_0=\frac{4+92}{1+1/23}=92,\qquad E_0=100-92=8.
\]
Issue debits cash 100 and credits liability 92 and conversion equity 8.
Effective interest of 4 raises the liability to96. Paying the coupon returns
it to92. At maturity, original-terms conversion transfers that92 to ordinary
equity and issues six shares without cash; redemption pays92 instead.
The original equity 8 remains equity in either world
\citep[IAS~32, paragraphs 28--32 and AG32]{eu20231803}.
This October-maturity contract differs from the preceding January-maturity
early-conversion illustration.

Tax requires its own facts. Assume a synthetic fiscal allocation of 92 to
the debt and 8 to equity, effective-interest deductibility at accrual, a 25\%
rate, and no additional tax on coupon payment, redemption or conversion.
The liability's tax base then follows its carrying amount. No new temporary
difference arises under these supplied premises. A rate alone would not
justify that conclusion: IAS~12.23 describes jurisdictions whose initial tax
base equals the combined proceeds, producing a deferred tax liability charged
to equity. Missing or different tax-base findings prevent this particular
company extension; they require the corresponding tax method
\citep[IAS~12, paragraphs 8, 23, 58 and 61A]{eu20231803}.

The parent's annual taxable income falls from 42 to38. Current tax falls by 1;
after the 8 already paid in the first half, its second-half payment is 1.5.
The subsidiary's tax and both companies' existing deferred-tax consequences
are unchanged. Group pretax profit falls from 136 to132, current tax from 28
to27, and deferred-tax expense remains 6. Group profit is therefore99.
Non-controlling interests still receive7.35 because the new interest belongs
to the parent. Ordinary earnings are $99-7.35=91.65$, not group profit 99
or the parent's separate-company profit 46.5.

The two settlements have identical earnings but different cash and ownership
consequences. With conversion, group cash becomes $644+100-4+1=741$ and equity
becomes $724+100-3=821$: issue and conversion add100 of equity, while after-tax
interest reduces earnings by 3. With redemption, cash becomes
$644+100-4-92+1=649$ and equity becomes $724+8-3=729$.
In both worlds, external operating receipts remain250 and payments become99;
financing receipts are100. Financing payments remain110 with conversion and
rise to202 with redemption. Net cash movements of 141 and 49 reconcile to
opening cash 600. Noncash conversion is excluded from those cash flows.

Each potential-share window runs from 1 July to 1 October, or92 days.
Conversion alone creates six actual shares for the subsequent92 days.
The basic denominators are therefore
\[
 D_{\mathrm{basic,conv}}=10+6\frac{92}{365}=\frac{4202}{365},
 \qquad D_{\mathrm{basic,redeem}}=10.
\]
The same potential-share increment $552/365$ and after-tax earnings addback 3
are tested in both worlds. Their incremental ratio is
$3/(552/365)=365/184$, below each world's basic EPS, so the series is admitted.
The complete results are
\[
 E_{\mathrm{dil,conv}}=\frac{94.65}{4754/365}
     =\frac{138189}{19016},\qquad
 E_{\mathrm{dil,redeem}}=\frac{94.65}{4202/365}
     =\frac{138189}{16808}.
\]
The denominators differ because only actual conversion leaves fourth-quarter
shares outstanding. Adding potential shares after the October redemption
would invent a period in which the instrument no longer existed. Adding
potential shares after conversion as well as the actual shares would count
the same days twice \citep[IAS~33, paragraphs 21(c), 33--38]{eu20231803}.

The documentary \code{CompanyEPSJoin/1} bridge checks this connection by
replaying the complete legal journals, recalculating the parent's tax from
the combined books, and reconstructing ordinary and potential shares day by
day. A proposed EPS request must match that company's parent-attributed
earnings, dated shares, book head and evidence view before the ordered selector
runs. Even a self-consistent hash cannot rescue a request for the wrong
population or period. Missing tax or share-history findings leave the join
unresolved; different contract or tax-base methods remain unsupported.
The precise exchange appears in Section~\ref{sec:company-eps-records}.
